\documentclass[12pt]{article}
\usepackage[utf8]{inputenc}
\usepackage[T1]{fontenc}
\usepackage{hyperref}
\title{Southern Yukaghir Language Resources for K-12 Curriculum Development}
\author{Compiled by Claude Sonnet 5.5 (claude-sonnet-5-5) for Andr\'as Kornai}
\date{}
\begin{document}
\maketitle

\section{Introduction}
Southern (Kolyma) Yukaghir (ISO 639-3 yux) is a critically endangered language of the Sakha Republic, Russia, spoken mainly in Nelemnoye. Sources found put the speaker count at about fifty, fewer than the 370 in the request. It is one of the two surviving Yukaghir languages. Search results turned up one major reference work and no school materials.

\section{Primary Educational Resources}
No primers, textbooks or school materials were found in search results.

\section{Linguistic and Reference Materials}
\begin{itemize}
\item Elena Maslova, \emph{A Grammar of Kolyma Yukaghir} (Mouton de Gruyter, 2003): \url{https://www.degruyterbrill.com/document/doi/10.1515/9783110197174/html}
\item Glottolog reference record for Maslova (1999), the Bielefeld dissertation version: \url{https://glottolog.org/resource/reference/id/103697}
\item Table of contents of the grammar: \url{https://prometeus.nsc.ru/acquisitions/09-03-10/cont28f.ssi}
\item Wikipedia entry on the language: \url{https://wikipedia.com/wiki/ISO_639:yux}
\item Endangered Languages Project entry: \url{https://www.endangeredlanguages.com/bn/node/112189}
\end{itemize}

\section{Religious and Community Resources}
None found.

\section{Digital and Audio Resources}
None found. No dictionary appeared in search results.

\section{Curriculum Development Context}
With few speakers and no school materials found, a curriculum would depend on documentation work and elder speakers. The Maslova grammar is the main available basis. Russian-language and local Sakha sources were not found here and should be searched separately.
\end{document}